\centerline{\bf \Huge Инвариант}


\begin{flushright}
        \scriptsize{Инвариант с рынка}
\end{flushright}

\problem{0.1} 
Можно ли доску 5 × 5 заполнить костяшками домино размером 1 × 2?

\problem{0.2} 
На вешалке висят 20 платков. 17 девочек по очереди подходят к вешалке, и каждая либо снимает, либо вешает ровно один платок. Может ли после ухода девочек на вешалке остаться 10 платков?

\problem{0.3}
В марсианском алфавите есть две буквы - У и Ы, причем если из любого слова выкинуть стоящие рядом буквы УЫ, то смысл слова не изменится. Точно также смысл не изменится при добавлении в любое место слова буквосочетания ЫУ или УУЫЫ. Верно ли, что слова ЫУЫУЫ и УЫУЫУ имеют одинаковый смысл?


\problem{1} 
Может ли шахматный слон за 100 000 ходов попасть с поля А1 на поле А8?


\problem{2} 
Дядька Черномор написал на листке бумаги число 20. Тридцать три богатыря передают листок друг другу, и каждый или прибавляет к числу, или отнимает от него единицу. Может ли в результате получиться число 10?

\problem{3}Из шахматной доски вырезали две противоположные угловые клетки. Докажите, что оставшуюся фигуру нельзя разре- зать на прямоугольники 1 × 2.



\problem{4} Набор чисел a, b, c каждую секунду заменяется на $a + b - c$, $a + c - b$, $b + c - a$. В начале имеется набор чисел 2015, 2016, 2017. Может ли через некоторое время получиться набор 2014, 2015, 2016?

\problem{5}
За один ход разрешается менять все знаки в одной строке или в одном столбце на противоположные. Можно
таким операциями привести таблицы, показанные на рисунке, к таблице из одних плюсов?

\begin{center}
\includegraphics[width=0.6\textwidth]{Летние Сборы/Сборы Элиста 2024/ВВ/инвариант_плюсы.png}
\end{center}


